\section{Demografische und nutzungsbezogene Merkmale der Teilnehmenden}
\label{sec:BB-03_demografie}

\begin{table}[H]
\centering
\caption{Demografische und nutzungsbezogene Merkmale der Teilnehmenden}
\label{tab:BB-demografie}
\begin{tabularx}{\textwidth}{l c l l l c}
\toprule
TN-ID & Alter & Geschlecht & Wikinutzung & Persona & Kompetenz \\
\midrule
TN-01 & 37 & weiblich & mehrmals pro Woche & PO Mitglied & 5 \\
TN-02 & 39 & männlich & einmal pro Woche & PO Mitglied & 4 \\
TN-03 & 40 & männlich & einmal pro Woche & PO Mitglied & 5 \\
TN-04 & 31 & weiblich & mehrmals im Monat & PO Mitglied & 4 \\
TN-05 & 52 & männlich & einmal im Monat & Keine PO Führungskraft & 4 \\
TN-06 & 23 & männlich & nie & PO Mitglied & 5 \\
TN-07 & 36 & männlich & einmal im Monat & PO Führungskraft & 5 \\
TN-08 & 23 & weiblich & nie & Kein PO Mitglied & 5 \\
TN-09 & 22 & weiblich & nie & Kein PO Mitglied & 5 \\
TN-10 & 45 & männlich & einmal pro Woche & PO Führungskraft & 5 \\
\bottomrule
\end{tabularx}

\vspace{0.5em}
\footnotesize
\textit{These zur Selbsteinschätzung der digitalen Kompetenz:} Ich fühle mich im Umgang mit digitalen Tools (z.\,B. Confluence, Intranet, Kollaborationstools) sicher.
\textit{Skala digitale Kompetenz:} 
1 = stimme gar nicht zu,\;
2 = stimme eher nicht zu,\;
3 = neutral,\;
4 = stimme eher zu,\;
5 = stimme überwiegend zu.\\
\end{table}
